\subsection*{NT-E13. The endpoint obstruction descends to potential CM curves}

Let \(E_1,E_2\) be elliptic curves over a number field \(k\)
with distinct imaginary quadratic potential-CM fields
\(M_1,M_2\). All endomorphisms need not be defined over \(k\).
Put \(F_i=kM_i\), \(F=F_1F_2\), and let \(S\) contain
the bad-reduction primes of both curves and the finite
primes ramified in \(F/k\). At any prime
\(\mathfrak p\notin S\), write \(q=N_k\mathfrak p\)
and let \(a_i\) be the two Frobenius traces.
Then
\[
       a_1^2=a_2^2\ne0
       \quad\Longrightarrow\quad
       \begin{cases}
       \mathfrak p\text{ splits completely in }F/k,\\
       q\text{ is a square and }a_1^2=4q .
       \end{cases}
                                                               \tag{E13.1}
\]
Consequently the simultaneous integral extreme
\[
              |a_1|=|a_2|=\lfloor2\sqrt q\rfloor
                                                               \tag{E13.2}
\]
is impossible at every absolute-degree-one prime outside
\(S\). Among all prime ideals outside \(S\) of norm at most
\(x\), its count is at most \([k:\mathbf Q]\sqrt x\).
These assertions are unconditional.
The nonzero condition in (E13.1) is essential: both curves
may have trace zero in nonsplit CM components.

\begin{proof}
For a characteristic-zero CM elliptic curve the field
over which all endomorphisms are defined is \(kM_i\).
One way to see the exact quadratic action is to let an
endomorphism act on the one-dimensional tangent space.
In characteristic zero this is a faithful map from the
CM field into \(\overline k\), compatible with Galois
action. The action on endomorphisms therefore factors
through the two automorphisms of \(M_i\), with kernel
precisely the group fixing \(kM_i\).
This is the same endomorphism-field premise used in
NT-E08 and NT-E10.

Fix an auxiliary prime \(\ell\) different from the
residue characteristic. Over an algebraic closure of
\(\mathbf Q_\ell\), the faithful CM action decomposes
the two-dimensional Tate module into its two distinct
one-dimensional eigenspaces. For an arithmetic
Frobenius at \(\mathfrak p\), conjugation of
endomorphisms agrees with its action on \(M_i\).
If the restriction to \(F_i/k\) is nontrivial, Frobenius
interchanges those two eigenspaces. Its matrix in that
decomposition has zero diagonal, hence trace zero.
The Frobenius trace is \(a_i\), since the curve has
good reduction and \(\mathfrak p\) is unramified in
the indicated quadratic extension.

Thus \(a_i\ne0\) forces trivial Frobenius in \(F_i/k\).
If \(a_1^2=a_2^2\ne0\), the Frobenius is trivial on
both \(F_i\), hence on their compositum \(F\).
This proves the splitting assertion in (E13.1).
Choose a prime \(\mathfrak P\) of \(F\) over \(\mathfrak p\).
It has the same residue field, norm \(q\), and both
traces \(a_i\). The two base-changed curves now have
all CM defined over \(F\), with the same distinct
fields \(M_1,M_2\). NT-E12 applies and forces square
norm and exact Hasse equality, proving (E13.1).

Since \(q\ge2\), its extremal integer
\(\lfloor2\sqrt q\rfloor\) is positive. Event (E13.2)
therefore satisfies the nonzero premise of (E13.1).
It cannot occur at norm a rational prime.
For a square norm \(q=p^{2f}\le x\), the underlying
rational prime is at most \(\sqrt x\), and at most
\([k:\mathbf Q]\) prime ideals of \(k\) lie above each
one. This gives the stated count bound.
\end{proof}

\paragraph{An explicit pair over \(\mathbf Q\).}
For
\[
        E_1:y^2=x^3-x,\qquad E_2:y^2=x^3+1,
                                                               \tag{E13.3}
\]
there is no rational prime \(p>3\) at which both
\(|a_p(E_i)|=\lfloor2\sqrt p\rfloor\).
Indeed the automorphism
\((x,y)\mapsto(-x,iy)\) gives CM by \(\mathbf Z[i]\)
on \(E_1\), while
\((x,y)\mapsto(\zeta_3x,y)\) gives CM by
\(\mathbf Z[\zeta_3]\) on \(E_2\).
These are the maximal orders in the distinct
quadratic fields \(\mathbf Q(i)\) and
\(\mathbf Q(\sqrt{-3})\); their displayed generators
are defined over those fields.
The two integral equations have discriminants
\(64\) and \(-432\), respectively, so both curves
have good reduction at every \(p>3\).
The compositum is \(\mathbf Q(\zeta_{12})\), ramified
only at \(2,3\). Every \(p>3\) is therefore within
the hypotheses just proved.


\paragraph{Earlier rational special case.}
Bucur--Fit\'e--Kedlaya, Remark~4.11, published p.~1743
(author v2 p.~22), already gives the obstruction
$a_p(E_1)=a_p(E_2)=\lfloor2\sqrt p\rfloor$ over
$\mathbf Q$ for distinct CM fields and $p>3$,
crediting Andrew Sutherland. The formulation here
retains arbitrary base fields, equal absolute traces,
the precise square-norm exception, and the quantitative
count of higher-degree primes. The rational positive-extreme
special case is not claimed as a new literature result.

\paragraph{Input, scope and geometric placement.}
The proof combines the exact CM component action with
the rational-Frobenius exception of NT-E12. It does
not presume that a nonsplit component has a continuous
circle distribution; its trace is zero instead.
It gives a concrete obstruction over \(\mathbf Q\)
at the total shrinking exponent \(1/2\), and remains
true after any further finite Frobenius restriction.
The relevant geometry is the component group acting
on the two CM eigenlines, followed on the split
component by the two distinct quadratic endomorphism
algebras. No GRH or Artin-holomorphy input occurs.
Full proof: \texttt{potential\_cm\_joint\_endpoint.tex},
NT-E13; run \texttt{GRH\_NIGHT/number\_theory}.
